\chapter{Hiring and Compensation}\label{ch:fm:hiring-and-compensation}

A trader paid fifteen per cent of a twenty-million-dollar year and a trader paid nothing after a two-million-dollar loss on the next desk may have
exactly the same skill. For a trader whose risk gives an annual P\&L volatility of \$10 million and whose true Sharpe ratio is the desk's
average of 0.8, the first year happens about one year in nine and the second one year in six. One year of P\&L measures luck far more than ability: on
the chapter's desk, 87\% of the differences in a year's formulaic pay between traders are luck. A pay rule decides which of the two the firm is
buying, and how long it takes to find out.

\section{The pipeline and the interview}

A trading firm hires from a funnel: applications, screens, tests, interviews, offers, acceptances. Two facts shape it. The base rate is low, so a
test with a modest false-positive rate still fills the later stages with people who will not be hired; and every stage's decision is a noisy
measurement, so the order of stages matters: cheap, valid tests first, expensive ones last.

\begin{definition}[Structured interview]\label{def:fm:hiring-and-compensation:structured}
A \emph{structured interview}\index{structured interview} asks every candidate for a role the same questions, chosen from an analysis of the job,
and scores the answers against anchored scales agreed in advance, with several interviewers scoring independently before they confer.
\end{definition}

The research on selection is unusually settled on one point. Schmidt and Hunter's review of 85 years of studies (1998) and its revision by
Sackett, Zhang, Berry and Lievens (2022), which found that earlier estimates had been overstated by 0.10 to 0.20 through over-corrections, agree
that \omterm{def:fm:hiring-and-compensation:structured}{structured interviews} predict job performance better than unstructured ones; in the revision they are the top-ranked selection procedure.
The interview questions of this series are written for that purpose: the same question, a known good answer, and what the interviewer is looking
for. One Quant Book 18 applies the same format to new questions.

\begin{method}[Designing an interview loop]\label{met:fm:hiring-and-compensation:loop}
\begin{enumerate}
\item Write the role's analysis: the three to five abilities that separate good performers from the rest.
\item For each ability choose a test or structured questions, with anchored scores.
\item Put cheap, valid tests first (a timed problem set, a coding test) and expensive ones (on-site loops) last.
\item Score independently before discussing; record scores and outcomes, and check a year later which scores predicted performance.
\end{enumerate}
\end{method}

\section{Pay structures: base, bonus and deferral}

Pay in trading has three parts: base salary, fixed for the year; a bonus, decided after it; and, above some level, a deferred part that vests
over later years. The balance between them is the firm's choice where no rule limits it and a regulator's where one does.

\begin{definition}[Bonus pool, formulaic payout, discretionary bonus]\label{def:fm:hiring-and-compensation:pool}
A \emph{bonus pool}\index{bonus pool} is the total variable pay a firm or desk awards for a year, set by a rule or a decision from its results.
A \emph{formulaic payout}\index{formulaic payout} pays a trader or team a fixed share of their own P\&L, usually net of costs and of losses
carried forward (chapter 3). A \emph{discretionary bonus}\index{discretionary bonus} is set by management's judgement of the individual's
contribution, within the pool.
\end{definition}

\begin{definition}[Deferred compensation, vesting schedule]\label{def:fm:hiring-and-compensation:deferral}
\emph{Deferred compensation}\index{deferred compensation} is variable pay awarded for a year but paid in later years, in cash, shares or fund
units. Its \emph{vesting schedule}\index{vesting schedule} sets when each part becomes the employee's: all at the end of a period (cliff) or in
equal parts each year (pro rata).
\end{definition}

\begin{definition}[Malus, clawback]\label{def:fm:hiring-and-compensation:malus}
\emph{Malus}\index{malus} is the reduction or cancellation of variable pay that has been awarded but not yet vested, for misconduct, a material
error or losses that come to light later. \emph{Clawback}\index{clawback} is the recovery of variable pay that has already been paid or has
vested.
\end{definition}

\begin{definition}[Material risk taker]\label{def:fm:hiring-and-compensation:mrt}
A \emph{material risk taker}\index{material risk taker} is an employee whose activities, under banking and investment-firm remuneration
rules, have a material impact on the institution's risk profile; their variable pay is subject to the rules on deferral, instruments, \omterm{def:fm:hiring-and-compensation:malus}{malus} and
\omterm{def:fm:hiring-and-compensation:malus}{clawback}.
\end{definition}

\begin{dated}{2026-09}{The pay rules}\label{dat:fm:hiring-and-compensation:rules}
\textbf{European Union} (Directive 2013/36/EU, article 94, as amended by Directive (EU) 2019/878): for \omterm{def:fm:hiring-and-compensation:mrt}{material risk takers} of banks, variable pay
may not exceed 100\% of fixed pay, or 200\% with shareholders' approval; at least half is paid in shares or equivalent instruments; at least 40\%
is deferred over four to five years, vesting no faster than pro rata (60\% for particularly high amounts); up to 100\% is subject to \omterm{def:fm:hiring-and-compensation:malus}{malus} or
\omterm{def:fm:hiring-and-compensation:malus}{clawback}; guaranteed variable pay is allowed only when hiring, for the first year. \textbf{United Kingdom}: the PRA and the FCA removed the ratio
limit (the ``bonus cap'') for banks, building societies and PRA-designated investment firms from 31 October 2023 (PRA PS9/23, FCA PS23/15), citing
its effect of raising fixed pay. \textbf{United States}: SEC Rule 10D-1 requires listed companies to recover incentive pay erroneously awarded to
executive officers in the three fiscal years before an accounting restatement. \textbf{High earners}: the European Banking Authority counted 2\,554
people earning more than \euro1 million at EU banks and investment firms in 2024, up from 2\,343 in 2023; in investment firms the number rose 30\%,
from 221 to 288.
\end{dated}

\section{Sizing a bonus pool}

A pool is sized top-down, then allocated. The common rules are a share of revenue (the \omterm{def:fm:the-economics-of-a-trading-firm:comp}{compensation ratio} of chapter 1, fixed in advance) or a
share of net profit after costs and a charge for the capital used (chapter 7). The second pays only for value added: a desk that earns less than its
cost of capital has no pool, whatever its revenue. The choice matters most in the years the firm least wants to pay: a revenue-based pool pays out
in a year the desk destroyed value on its capital.

\begin{example}[Two pools]\label{ex:fm:hiring-and-compensation:pools}
A desk earns revenue of 100 with costs of 60 and a \omterm{def:fm:limits-and-capital-allocation:raroc}{capital charge} of 10. A pool of 20\% of profit after the charge is $0.2\times(100-60-10)=6$. A
pool of 6\% of revenue is also 6 this year. In a year of revenue 70, the first pool is zero and the second is 4.2.
\end{example}

\section{Formulaic against discretionary: paying for skill, not luck}

A trader's P\&L in a year is skill plus luck. If skill across traders has a standard deviation $s$ (in P\&L per year) and luck, the year's noise, a
standard deviation $n$, then after $T$ years the average P\&L correlates with skill at $\sqrt{s^2/(s^2+n^2/T)}$.

\begin{proposition}[How long pay takes to find skill]\label{prop:fm:hiring-and-compensation:years}
Let trader $i$'s annual P\&L be $\mu_i+\varepsilon_{i,t}$, with $\mu_i$ drawn with variance $s^2$ and $\varepsilon_{i,t}$ independent with
variance $n^2$. The share of the cross-sectional variance of one year's P\&L due to luck is $n^2/(s^2+n^2)$, and the correlation of $T$-year average
P\&L with $\mu_i$ reaches $\rho$ after
\[
T=\frac{n^2}{s^2\,(1/\rho^2-1)}\ \text{years}.
\]
\end{proposition}

\begin{proof}
The average is $\mu_i+\bar\varepsilon_i$ with $\mathrm{Var}(\bar\varepsilon_i)=n^2/T$; its covariance with $\mu_i$ is $s^2$, so its correlation
is $s/\sqrt{s^2+n^2/T}$. Setting it to $\rho$ and solving gives $T$.
\end{proof}

On the chapter's desk (illustrative: twenty traders, annual P\&L volatility \$10 million each, true Sharpe ratios spread with a standard deviation
of 0.4, so $s=4$ and $n=10$), luck is 86\% of the variance of one year's P\&L, and it takes 11.1 years of P\&L for it to correlate with skill at 0.8.
A \omterm{def:fm:hiring-and-compensation:pool}{formulaic payout} of 15\% of positive P\&L inherits both: 87\% of the variance of a year's pay is luck, and its correlation with skill is 0.35 after
one year, 0.64 after five and 0.80 after about twelve (\cref{fig:fm:hiring-and-compensation:years}).

\begin{omfigure}
\begin{tikzpicture}
\begin{axis}[width=10.5cm,height=5.2cm,xlabel={\small years of P\&L},ylabel={\small correlation of pay with skill},xmin=1,xmax=15,ymin=0,ymax=1,
  tick label style={font=\footnotesize},grid=major,grid style={gray!20},table/col sep=comma,
  legend cell align=left,legend style={font=\scriptsize,at={(0.5,-0.3)},anchor=north,/tikz/every even column/.append style={column sep=8pt}},legend columns=2]
\addplot[omDef,thick,mark=*,mark size=1.2pt] table[x=years,y=simulated]{figdata/desk/10-hiring-and-compensation/corr.csv};
\addplot[omThm,thick,dashed] table[x=years,y=linear]{figdata/desk/10-hiring-and-compensation/corr.csv};
\draw[gray,dotted] (axis cs:1,0.8) -- (axis cs:15,0.8);
\legend{cumulative formulaic pay (simulated),{average P\&L (\cref{prop:fm:hiring-and-compensation:years})}}
\end{axis}
\end{tikzpicture}

\omcaption{How long pay takes to reflect skill: the correlation across twenty traders between true Sharpe ratio and cumulative pay under a 15\%
\omterm{def:fm:hiring-and-compensation:pool}{formulaic payout} (400 simulated desks), and between skill and average P\&L in closed form. Skill spread 0.4 in Sharpe ratio, annual volatility
\$10 million. Data: \texttt{fm\_pay.corr\_by\_years}.}\label{fig:fm:hiring-and-compensation:years}
\end{omfigure}

A discretionary rule can do better only if it uses information the P\&L does not contain: the quality of the trader's decisions, their risk-taking,
their contribution to others' P\&L. Judgement applied to the P\&L alone is at best a shrinkage estimate of skill
(\cref{lst:fm:hiring-and-compensation:shrink}), and on the chapter's desk it correlates with skill at 0.37 after one year against the formulaic 0.35,
and slightly less after five (0.61 against 0.64), because shrinkage is proportional and the pool it shares moves with the desk's luck. A
risk-adjusted payout, charging each trader's capital, correlates at 0.34 and 0.63: it pays for value added, not for skill detection.

\omcode{code/firm/bonuspool/firm_bonuspool.py}{44}{56}{A discretionary allocation that uses only P\&L is a shrinkage of each trader's record towards
the desk's prior.}\label{lst:fm:hiring-and-compensation:shrink}

\begin{omfigure}
\begin{tikzpicture}
\begin{axis}[width=10.5cm,height=5.2cm,xlabel={\small true Sharpe ratio},ylabel={\small five-year pay (\$ million)},xmin=-0.4,xmax=2.0,ymin=0,
  tick label style={font=\footnotesize},grid=major,grid style={gray!20},table/col sep=comma]
\addplot[only marks,mark=*,mark size=1.4pt,omDef] table[x=sharpe,y=pay]{figdata/desk/10-hiring-and-compensation/scatter.csv};
\end{axis}
\end{tikzpicture}

\omcaption{Five years of formulaic pay (15\% of positive annual P\&L) against the trader's true Sharpe ratio, for 100 traders on five simulated
desks. Pay rises with skill, and traders of the same skill are paid anything from almost nothing to twice the average. Data:
\texttt{fm\_pay.one\_desk}.}\label{fig:fm:hiring-and-compensation:scatter}
\end{omfigure}

\begin{remark}[Why formulaic pay persists]\label{rem:fm:hiring-and-compensation:why}
\omterm{def:fm:hiring-and-compensation:pool}{Formulaic payouts} are common where the firm cannot or will not judge skill, and where the trader carries their own drawdown (chapter 3's
platforms): the formula is transparent and hard to dispute, and it attracts people who believe in their own edge. Its cost is that it pays luck, and
that it makes a trader's pay convex in their P\&L, which rewards risk unless limits and deferral cap it.
\end{remark}

\section{Deferral, clawback and retention}

Deferral does three jobs. It exposes pay to the losses that appear after the year (through \omterm{def:fm:hiring-and-compensation:malus}{malus}); it aligns pay with the firm's longer-term
results; and it builds a balance that the employee loses by leaving. The chapter's plan pays 60\% of each award in cash and vests the rest in four
equal annual parts (\cref{fig:fm:hiring-and-compensation:vesting}). A trader who receives the same award every year carries, at each year end, an
unvested balance equal to one full year's award: $0.4\times(1+\tfrac34+\tfrac12+\tfrac14)=1$.

\begin{definition}[Deferral buyout]\label{def:fm:hiring-and-compensation:buyout}
A \emph{deferral buyout}\index{deferral buyout} is the payment a hiring firm makes to a new employee to replace the deferred pay forfeited by
leaving their previous employer, usually in the new firm's own deferred instruments and on a similar schedule.
\end{definition}

\omcode{code/firm/bonuspool/firm_bonuspool.py}{59}{79}{A deferral plan, the payments of one award, and the unvested balance at a year end.}\label{lst:fm:hiring-and-compensation:defer}

\begin{omfigure}
\begin{tikzpicture}
\begin{axis}[width=10.5cm,height=5cm,ybar stacked,bar width=14pt,xlabel={\small year},ylabel={\small paid (units of one award)},xmin=-0.6,xmax=7.6,
  ymin=0,ymax=1.1,xtick={0,1,2,3,4,5,6,7},tick label style={font=\footnotesize},grid=major,grid style={gray!20},table/col sep=comma,
  legend cell align=left,legend style={font=\scriptsize,at={(0.5,-0.3)},anchor=north,/tikz/every even column/.append style={column sep=6pt}},legend columns=5]
\addplot[fill=omDef!55,draw=omDef] table[x=year,y=c0]{figdata/desk/10-hiring-and-compensation/vesting.csv};
\addplot[fill=omProp!55,draw=omProp] table[x=year,y=c1]{figdata/desk/10-hiring-and-compensation/vesting.csv};
\addplot[fill=omThm!50,draw=omThm] table[x=year,y=c2]{figdata/desk/10-hiring-and-compensation/vesting.csv};
\addplot[fill=gray!40,draw=gray] table[x=year,y=c3]{figdata/desk/10-hiring-and-compensation/vesting.csv};
\addplot[fill=omDef!20,draw=omDef,dashed] table[x=year,y=c4]{figdata/desk/10-hiring-and-compensation/vesting.csv};
\legend{award 0,award 1,award 2,award 3,award 4}
\end{axis}
\end{tikzpicture}

\omcaption{Payments from five equal annual awards under a plan that pays 60\% in cash and vests 40\% pro rata over four years. From year 4 the
employee receives a full award a year, and always has one full award unvested. Data: \texttt{firm.bonuspool.schedule}.}\label{fig:fm:hiring-and-compensation:vesting}
\end{omfigure}

The balance is the retention device and the price of a move. A competitor that hires the trader must buy out a year's award in deferred
instruments of its own, and the trader carries the firm's \omterm{def:fm:hiring-and-compensation:malus}{malus} risk for the years the balance vests. \omterm{def:fm:hiring-and-compensation:malus}{Clawback} reaches further than \omterm{def:fm:hiring-and-compensation:malus}{malus}, to pay
already received, and is harder to enforce: the US rule for listed companies' executives, the EU banking rules and most private firms' plans
all rely mainly on \omterm{def:fm:hiring-and-compensation:malus}{malus}, applied to what is still unvested.

\section{Tutorial: twenty traders and three pay rules}\label{tut:fm:hiring-and-compensation:tut}

\begin{tutorial}
\textbf{Goal.} Measure how well each pay rule tracks skill, and the retention a deferral plan creates. \textbf{End state:}
\cref{fig:fm:hiring-and-compensation:years} and the correlations of the text.
\begin{tutsteps}
\item \textbf{The desk.} \texttt{fm\_pay.desk(rng)} draws twenty traders' true Sharpe ratios and five years of P\&L.
\item \textbf{The rules.} \texttt{fm\_pay.pay(pnl, rule)} applies \texttt{firm.bonuspool.formulaic}, \texttt{risk\_adjusted} and
\texttt{discretionary}.
\item \textbf{The correlations.} \texttt{fm\_pay.correlations()} averages, over 2\,000 desks, the correlation of one and five years' pay with skill;
\texttt{luck\_share\_formulaic()} computes the luck share of a year's pay.
\item \textbf{Deferral.} \texttt{firm.bonuspool.unvested([1.0]*6, DeferralPlan(0.6, 4), 5)} gives the steady unvested balance, one award.
\end{tutsteps}
\end{tutorial}

\textbf{What to change next.} Double the skill spread and watch the years needed fall by four; give the discretionary rule a noisy observation of
each trader's decision quality alongside the P\&L.

\section{Build: the bonus pool}\label{bld:fm:hiring-and-compensation:bonuspool}

\begin{build}
\bfield{Purpose} The desk's pay arithmetic: pool sizing, allocation rules, deferral and its balances, and how much of pay is luck.
\bfield{Interface} \texttt{firm.bonuspool}: \texttt{pool\_from\_profit}, \texttt{pool\_from\_revenue}; \texttt{formulaic}, \texttt{risk\_adjusted},
\texttt{shrunk\_skill}, \texttt{discretionary}; \texttt{DeferralPlan(cash\_share, years)}, \texttt{schedule}, \texttt{unvested}, \texttt{malus};
\texttt{luck\_share\_linear}, \texttt{years\_for\_correlation}.
\bfield{Rules} A pool is never negative; a discretionary allocation adds up to the pool; deferred awards vest pro rata; \omterm{def:fm:hiring-and-compensation:malus}{malus} acts only on the
unvested balance.
\bfield{Acceptance tests} \texttt{code/firm/bonuspool/tests/}: pools and rules on hand numbers; the shrinkage weight; the schedule and the steady
unvested balance; the closed forms of \cref{prop:fm:hiring-and-compensation:years}.
\bfield{Stretch} Pay in fund units whose value tracks the desk's later P\&L; a guaranteed first-year award and its cost; the pool allocated by
Shapley credit (chapter 9).
\end{build}

\begin{omsources}
\item F.~L.~Schmidt and J.~E.~Hunter, ``The validity and utility of selection methods in personnel psychology'', \emph{Psychological Bulletin}
124(2), 1998; P.~R.~Sackett, C.~Zhang, C.~M.~Berry and F.~Lievens, ``Revisiting meta-analytic estimates of validity in personnel selection'',
\emph{Journal of Applied Psychology} 107(11), 2022.
\item Directive 2013/36/EU, article 94, and Directive (EU) 2019/878; PRA PS9/23 and FCA PS23/15 (2023); 17 CFR 240.10D-1.
\item European Banking Authority, press release on high earners in 2024, 16 April 2026.
\end{omsources}

\section{Exercises}

\begin{exercise}[$\star$]\label{exo:fm:hiring-and-compensation:1}
A trader's expected annual P\&L is \$8 million with a volatility of \$10 million. What are the probabilities of a year of \$20 million or more and
of a loss of \$2 million or more?
\end{exercise}

\begin{exercise}[$\star$]\label{exo:fm:hiring-and-compensation:2}
Compute the two pools of \cref{ex:fm:hiring-and-compensation:pools} in a year of revenue 120.
\end{exercise}

\begin{exercise}[$\star$]\label{exo:fm:hiring-and-compensation:3}
A plan pays 60\% of each award in cash and vests the rest over four years pro rata. What is paid, in total, in the fifth year of a trader who has
received an award of 1 every year?
\end{exercise}

\begin{exercise}[$\star\star$]\label{exo:fm:hiring-and-compensation:4}
With skill spread $s=4$ and noise $n=10$, compute the luck share of one year's P\&L and the years needed for a correlation of 0.8. Redo it for
$s=8$.
\end{exercise}

\begin{exercise}[$\star\star$]\label{exo:fm:hiring-and-compensation:5}
Under the EU rules, a \omterm{def:fm:hiring-and-compensation:mrt}{material risk taker}'s variable pay is 150\% of fixed pay. What approval does that need, and what is the least that must be
deferred and paid in instruments?
\end{exercise}

\begin{exercise}[$\star\star$]\label{exo:fm:hiring-and-compensation:6}
Why does a \omterm{def:fm:hiring-and-compensation:pool}{formulaic payout} reward risk-taking, and which two tools in the chapter limit that?
\end{exercise}

\begin{exercise}[$\star\star\star$]\label{exo:fm:hiring-and-compensation:7}
\emph{Coding.} Rerun \texttt{fm\_pay.correlations(n\_desks=500, years=10)}. What is the ten-year correlation of formulaic pay with skill?
\end{exercise}

\begin{exercise}[$\star\star\star$]\label{exo:fm:hiring-and-compensation:8}
\emph{Find the flaw.} ``Our best trader made \$25 million last year, three times the desk average. Offer him a guaranteed \$4 million to stay.''
\end{exercise}

\section{Problem: Paid for Luck}

\begin{problem}[{Weekend problem --- paid for luck}]\label{pb:fm:hiring-and-compensation:1}
A head of desk must choose how to pay twenty traders, and must decide whether to match a competitor's offer to one of them.

\medskip\noindent\textbf{Part I --- Hiring.}
\begin{enumerate}
\item Define a \omterm{def:fm:hiring-and-compensation:structured}{structured interview} and say what the research finds about it.
\item Design a four-stage hiring loop for a trader and justify its order.
\item Why does a low base rate make early screening matter?
\end{enumerate}

\medskip\noindent\textbf{Part II --- The pool.}
\begin{enumerate}[resume]
\item Define a \omterm{def:fm:hiring-and-compensation:pool}{bonus pool}, a \omterm{def:fm:hiring-and-compensation:pool}{formulaic payout} and a \omterm{def:fm:hiring-and-compensation:pool}{discretionary bonus}.
\item Size a pool as 20\% of profit after costs and a \omterm{def:fm:limits-and-capital-allocation:raroc}{capital charge} for revenue of 100, costs of 60 and a charge of 10; and in a year of revenue 70.
\item Why does a pool on profit after a \omterm{def:fm:limits-and-capital-allocation:raroc}{capital charge} differ from one on revenue in a bad year?
\end{enumerate}

\medskip\noindent\textbf{Part III --- Skill and luck.}
\begin{enumerate}[resume]
\item State and prove \cref{prop:fm:hiring-and-compensation:years}.
\item Give the luck share of one year's P\&L and the years to a correlation of 0.8 on the chapter's desk.
\item Give the luck share of one year's formulaic pay.
\item Give each rule's correlation with skill after one and five years.
\item Why does discretion on the P\&L alone not beat the formula, and what would let it?
\item What is the ten-year correlation of formulaic pay with skill?
\end{enumerate}

\medskip\noindent\textbf{Part IV --- Deferral and the offer.}
\begin{enumerate}[resume]
\item Define \omterm{def:fm:hiring-and-compensation:deferral}{deferred compensation}, a \omterm{def:fm:hiring-and-compensation:deferral}{vesting schedule}, \omterm{def:fm:hiring-and-compensation:malus}{malus}, \omterm{def:fm:hiring-and-compensation:malus}{clawback} and a \omterm{def:fm:hiring-and-compensation:buyout}{deferral buyout}.
\item Give the steady unvested balance of the chapter's plan.
\item What must a competitor pay to hire a trader on that plan, and in what form?
\item State the EU limits on the ratio, deferral and instruments, and the UK change of 2023.
\item Why is \omterm{def:fm:hiring-and-compensation:malus}{malus} easier to apply than \omterm{def:fm:hiring-and-compensation:malus}{clawback}?
\item Should the head of desk match the offer to the trader who made \$25 million?
\item State the \emph{named result}: the luck share of a one-year \omterm{def:fm:hiring-and-compensation:pool}{formulaic payout}, and the years of P\&L after which pay tracks skill with a
correlation of 0.8.
\item In two sentences, write the desk's pay policy.
\end{enumerate}
\end{problem}

\section{Interview questions}

\begin{interviewq}[$\star$ \iqroles{trader}]\label{iq:fm:hiring-and-compensation:1}
Your firm pays you 15\% of your P\&L. Why might it still defer part of it?
\end{interviewq}

\begin{interviewq}[$\star$ \iqroles{trader, researcher}]\label{iq:fm:hiring-and-compensation:2}
What is the difference between \omterm{def:fm:hiring-and-compensation:malus}{malus} and \omterm{def:fm:hiring-and-compensation:malus}{clawback}?
\end{interviewq}

\begin{interviewq}[$\star\star$ \iqroles{researcher}]\label{iq:fm:hiring-and-compensation:3}
Traders' true Sharpe ratios differ with a standard deviation of 0.4 and each year's P\&L has a Sharpe-unit noise of 1. How many years until P\&L
ranks traders reliably?
\end{interviewq}

\begin{interviewq}[$\star\star$ \iqroles{trader}]\label{iq:fm:hiring-and-compensation:4}
A competitor offers you a higher \omterm{def:fm:the-multi-manager-platform:payout}{payout rate} but no \omterm{def:fm:hiring-and-compensation:buyout}{deferral buyout}. How do you compare the offers?
\end{interviewq}

\begin{interviewq}[$\star\star$ \iqroles{risk}]\label{iq:fm:hiring-and-compensation:5}
Why does a \omterm{def:fm:hiring-and-compensation:pool}{formulaic payout} on P\&L create an incentive to take risk, and how would you control it?
\end{interviewq}

\begin{interviewq}[$\star\star\star$ \iqroles{researcher, risk}]\label{iq:fm:hiring-and-compensation:6}
How would you design a pay rule that rewards skill faster than P\&L alone can reveal it?
\end{interviewq}
